\scanpage{5}\mspage{3}

\deleted{\unsure{$E$ est à degrés pairs}, et que} l'on voudra que
\[
 \overline{\left|\ \pi_i=0\quad\text{si }i\notin[0,2n]\ \right|}\tag*{(3.0)}
\]
intervalle \inserted{(\unsure{pair|fini})} fixé. \deleted{\unsure{Suppo|Supp}} On s'occupera ici de l'\uline{unicité}. \deleted{Supposons établie l'unicité}

Notons d'abord que
\[
 \unsure{E|\mathcal E}^0\ni u\iff (u\pi_i=\pi_i u\qquad\forall i)
\]
donc $E^0$ est le commutant de la $k$-algèbre engendrée par les $\pi_i$. Donc les $\pi_i$ sont dans le centre $\unsure{Z|\mathfrak Z|\mathfrak z}(E^0)$ de $E^0$, et ce sont des projecteurs \inserted{mutuellement} orthogonaux de somme 1 de cet anneau commutatif. Supposons déterminé déjà les $\pi_\alpha$ pour $\alpha\notin[i,2n-i]$, \inserted{$i<n$,} et déterminons $\pi_i$ et $\pi_{2n-i}$ [i.e $\pi_\alpha$ pour $\alpha\notin[i+1,2n-i-1]$. \inserted{C'est clair si $i=n-1$, car on connaît tous $\pi_\alpha$ pour $\alpha\ne n$, d'où aussi $\pi_n=1-\sum_{\alpha\ne n}\pi_\alpha$. Reste le cas $i<n-1$} On connaît donc
\[
 \begin{aligned}
 \varphi_i&=\deleted{\unsure{(1-\sum\pi_\alpha)}}(1-\sum_{\alpha<i\deleted{\unsure{\alpha}}}\pi_i)=\sum_{\alpha\ge i}\pi_\alpha,\quad\text{et}\quad
 \psi_{2n-i}=(1-\sum_{\alpha>2n-i}\pi_i)=\sum_{\alpha\le 2n-i}\pi_{\unsure{\alpha'|\alpha}}
 \end{aligned}
\]
donc aussi la composée
\[
 \unsure{\text{Or}}\quad \psi_{2n-i}\deleted{\unsure{\varphi_i}}E^{2n-2i}\varphi_i
 =\sum_\alpha\psi_{2n-i}\pi_{\alpha+(2n-2i)}E\pi_\alpha\varphi_i)
\]
Or
\[
 \pi_\alpha\varphi_i=\begin{cases}0&\text{si }\alpha<i\\ \pi_\alpha&\text{si }\alpha\ge i\end{cases}
 \quad\text{et}\quad
 \psi_{2n-i}\pi_{\alpha+(2n-2i)}
 =\begin{cases}
 0&\text{si }\inserted{\alpha+(2n-2i)>2n-i}\ \text{i.e. }\alpha>i\\
 \pi_{\alpha+(2n-2i)}&\text{si }\alpha\le i
 \end{cases}
\]
donc le terme d'indice $\alpha$ de la dernière $\sum$ est nul sauf si $\alpha=i$, auquel cas il est égal à
\[
 \pi_{\unsure{\alpha|i}+(2n-2i)}E\pi_i=\pi_{2n-i}E\pi_i=E_{2n-i,i}\,.
\]
Donc la connaissance des $\pi_\alpha$ pour $\alpha\notin[i,2n-i]$ implique la connaissance de $E^{2n-i}_{i}\deleted{\unsure{,i}}$. \deleted{\unsure{Considér|considéré}} La condition qui va impliquer l'unicité \deleted{\unsure{est}} de $u_i$,
% archive sheet number: 4
